\section{Related Work}
\paragraph{Conformal prediction for trajectories}
Split CP turns any predictor into one with marginal coverage guarantees \cite{vovk2005algorithmic,angelopoulos2023gentle} and
underlies conformal safe planning around predicted agents \cite{lindemann2023safe}. Multi-step and multi-modal trajectories have been
handled with copula-based intervals \cite{sun2024copula}, scenario-aware intervals in a Frenet frame on nuPlan
\cite{shu2025scenario} and multi-view calibration that fuses vehicle- and infrastructure-side observations \cite{chen2024v2i}; a
2025 survey catalogues these directions and the open gaps around uncertainty and generalisation \cite{abdelmadjid2025survey}, and a
dataset-design study finds that which data the model sees matters more than feature engineering \cite{demmler2025dataset}, a finding
Sec.~\ref{sec:where} echoes from the calibration side. Calibrating the predictive \emph{distribution} itself has also been
proposed \cite{pouria2026gaussian}. Locally weighted scores that divide the residual by an estimated spread
\cite{lei2018distribution} are the classical route to adaptive regions; our normalised score uses the forecaster's own predicted
scale in that role.

\paragraph{Safe planning and robot safety with conformal guarantees}
Beyond trajectory regions, conformal prediction has certified learned perception for navigation \cite{mei2024perceive}, adapted a
crowd-navigation policy's safety margin to shift \cite{yao2025crowdnav}, learned safety regions from sparse human feedback
\cite{feldman2025robotsafety}, certified learned motion primitives \cite{binny2026stable} and connected learning-enabled control to
formal verification \cite{lindemann2024formal}. Of these, \cite{mirzaeedodangeh2025interactive} transfers a safety guarantee across
policy updates analytically and \cite{binny2025moved} treats the shift the ego agent induces in others. All assume that calibration
and deployment \emph{domains} match; we study what happens when the domain itself changes.

\paragraph{Conformal prediction under shift}
Weighted CP handles covariate shift given the likelihood ratio \cite{tibshirani2019covariate}, and the coverage loss
under general shift can be bounded by a distance between calibration and test distributions \cite{barber2023beyond}.
Recent work replaces that bound with an explicit ambiguity-set radius under L\'{e}vy--Prokhorov perturbations
\cite{aolaritei2025levy}, regularises the representation with a Wasserstein penalty to shrink the achievable coverage
gap \cite{xu2025wasserstein}, or uses optimal transport with unlabelled target data to both estimate and mitigate the
loss in coverage \cite{correia2025nonexchangeable}. That optimal-transport method is the closest label-free competitor,
but its bracketing of the unlabelled test scores is built from the minimum and maximum score over all candidate labels,
for scores bounded in $[0,1]$, and is evaluated on classification and a one-dimensional binned regression task. For our
unbounded, metre-scale score over multi-modal trajectories that minimum is zero and the maximum is unbounded, and we
found no published construction for structured outputs, so we do not run it as a baseline; a trajectory-specific
bracketing is a natural extension. Robust conformal prediction instead guards against every test distribution inside
an $f$-divergence ball around the training distribution \cite{cauchois2024robust}; it restricts the shift class by
construction, as Proposition~\ref{prop:ident} says a useful label-free method must, but needs a radius chosen or
estimated for the shift, and we leave it as a baseline for future work. A parallel line gives sample-complexity lower bounds for when a selective predictor can certify a risk floor under covariate
shift at all \cite{liu2026certify}, a negative result in the spirit of our Proposition~\ref{prop:ident}. Online methods adapt
the threshold from realised errors: adaptive conformal inference \cite{gibbs2021adaptive}, conformal PID control
\cite{angelopoulos2023pid}, adaptive scores with a replay buffer \cite{tumu2026adaptnc}, and variants that revise past
predictions \cite{jun2026retrospective}, play a universal-portfolio game \cite{liu2026portfolio} or give training-conditional
regret bounds \cite{liang2026optimalregret}. The closest trajectory-specific instance is CUQDS \cite{huang2025cuqds}, which pairs a
Gaussian-process uncertainty module with online conformal PID control; like the others it needs a stream of ground-truth labels
at deployment, so it answers a different question and is not a baseline for our zero-label or fixed-labelled-batch settings.
Our ACI comparison (Sec.~\ref{sec:labelfree}) covers the same category without CUQDS's architecture-specific module. Other work
conditions on an observable nuisance parameter with a diffusion model \cite{rahaman2026shift} or treats the shift the ego vehicle
induces in others \cite{binny2025moved} (preprints). We study a different point: no feedback, no generative model, and three
real driving datasets.

\paragraph{Grouped, hierarchical and conditional coverage}
When observations come in groups, exchangeability across observations fails, and hierarchical conformal methods restore validity
by pooling or subsampling at the group level \cite{dunn2023hierarchical,lee2023hierarchical}, with extensions for hierarchical
outputs \cite{principato2024hierarchicaldata}, known group-level shift \cite{bhattacharyya2024groupweighted} and partial test-group
data \cite{mallick2026generalizedhcp}. The conditional-coverage literature asks for validity on subpopulations rather than on
average \cite{min2026unifiedtheory,zhou2026cpassessment,braun2025diagnostics,bairaktari2025kandinsky}. Risk-controlling procedures
\cite{bates2021rcps} and Learn-then-Test \cite{angelopoulos2025ltt} give finite-sample control of a monotone risk by testing a
fixed ordered set of thresholds, and betting-based $p$-values \cite{waudbysmith2024betting} make such tests far tighter than
Hoeffding-type bounds when losses are small. We combine the two for a setting those works do not address: dataset shift in
trajectory prediction, where the labelling unit is the \emph{scene} and the controlled risk is the scene-averaged miscoverage
of a fixed, source-calibrated forecaster. Valid $p$-values for bounded losses need care \cite{alvarez2024pvalue}; our own
implementation needed a corrected bet cap (Sec.~\ref{sec:scene}).

\paragraph{Risk control and betting, more generally}
Risk control has been extended to tradeoffs among several risks \cite{nguyen2024tradeoffs}, semi-supervised tuning
\cite{einbinder2024semisupervised} and post-processing of decision policies \cite{joshi2026riskcontrolled}; e-values give
anytime-valid alternatives to the $p$-values we use \cite{gauthier2025evalues}, and weighted martingales monitor
post-deployment shift \cite{prinster2025watch}. We use betting for a different purpose: certifying a \emph{scene-averaged}
risk from whole labelled scenes.

\paragraph{Cross-dataset trajectory prediction and perception under domain shift}
UniTraj harmonises several driving datasets and reports that accuracy drops when models are transferred between them
\cite{feng2024unitraj}; that study is about point accuracy, not about the validity of an uncertainty guarantee.
Beyond trajectories, \cite{ovadia2019can} showed that calibration of deep models degrades under dataset shift; our
question is the conformal analogue for structured, multi-modal forecasts. The same gap is documented on the perception
side: localisation uncertainty is miscalibrated under novel weather \cite{wu2025visuallocalization}, collaborative
perception generalises poorly across sensor rigs \cite{wang2025collaborative}, and cross-sensor pretraining is used to
close a domain gap that would otherwise need target labels \cite{kong2025largead,qiu2026crosssensor}. None of this asks
whether a \emph{conformal} guarantee survives the transfer; this paper measures, diagnoses and partially repairs that.
